% Standalone LaTeX table generated by scripts/01_make_tables.py
% Required packages: booktabs, array
% Suggested text for paper: The simulation uses dimensionless baseline values and ranges; these are modelling controls, not calibrated estimates.
\documentclass{article}
\usepackage[margin=0.8in]{geometry}
\usepackage{booktabs}
\usepackage{array}
\begin{document}
\begin{table}[ht]
\centering
\caption{Dimensionless baseline values and simulation ranges.}
\label{tab:parameter-ranges-v1}
\begin{tabular}{p{0.20\linewidth}p{0.20\linewidth}p{0.20\linewidth}p{0.20\linewidth}p{0.20\linewidth}}
\toprule
parameter & baseline & range\_min & range\_max & role \\
\midrule
gamma & 0.55 & 0.30 & 0.80 & long-memory exponent; avoid endpoints 0 and 1 \\
Delta & 1.0 & 0.5 & 2.0 & operational coarse-graining scale \\
A\_eff0 & 1.0 & 0.5 & 2.0 & effective signed-forcing intensity scale \\
L0 & 1.0 & 0.75 & 1.5 & baseline liquidity slope \\
z\_low & 0.05 & 0.01 & 0.10 & low-z response-regime marker \\
z\_high0 & 10.0 & 5.0 & 25.0 & high-z response-regime marker \\
alpha0 & 1.0 & 0.5 & 2.0 & baseline deterministic activity \\
a\_short & 0.25 & 0.0 & 0.75 & short-maturity activity lift \\
tau\_alpha & 0.50 & 0.25 & 1.0 & activity decay scale in years \\
L\_skew & 0.80 & 0.0 & 1.5 & lower liquidity on downside when positive \\
L\_smile & 1.50 & 0.0 & 3.0 & thinner liquidity away from center \\
nu\_skew & 1.00 & 0.0 & 2.0 & weaker downside resilience when positive \\
A\_stress & 0.75 & 0.0 & 1.5 & stronger downside forcing in combined profile \\
L\_floor & 0.20 & 0.10 & 0.50 & positive floor for liquidity profile \\
\bottomrule
\end{tabular}
\end{table}
\end{document}
